\section{Systematic Uncertainties}
\label{sec:syst}
\AnalysisTBD{complete systematics table.  The components identified so far
are: MC statistics; generic-$B\bar{B}$ composition (Sec.~\ref{sec:syst:bbbar});
fake-$D$ extrapolation; continuum normalisation; $D^{**}\ell\nu$ and gap-mode
modelling; $D^{*}\tau\nu$ and $D^{**}\tau\nu$ feed-down; form factors;
tracking and PID efficiency; fit bias and finite-template effects}

\subsection{Generic-$B\bar{B}$ composition}
\label{sec:syst:bbbar}
The treatment is to propagate the covariance of the tuned family weights
(Sec.~\ref{sec:bbbar}) into the fit as \texttt{normsys} constraints on the
generic-$B\bar{B}$ normalisations, replacing the free normalisation factors
of the current workspace (Sec.~\ref{sec:fit:modifiers}).
Because the family weights are strongly correlated, assigning one independent
constrained parameter per family would discard that correlation.
The intended construction is instead to diagonalise the covariance matrix and
use its eigenvectors as a small set of uncorrelated nuisance parameters, each
varying all five family weights coherently by
$\sqrt{\lambda_{i}}\, \hat{v}_{i}$.
\InProgress{this cannot be applied to the currently stored fit.  Two
obstacles, both consequences of the minimum sitting on a parameter bound
(Sec.~\ref{sec:bbbar:results}): the stored HESSE errors and the stored
covariance matrix disagree by up to three orders of magnitude for the two
at-limit parameters, and the covariance is near-singular (condition number
$\sim 10^{8}$) with two eigenvalues consistent with zero, so the construction
would assign essentially no uncertainty to two entire hadronic families.
The method is sound and the tooling exists
(\texttt{scripts/bbbar\_eigen\_systematics.py}); it requires a fit with an
interior minimum.  Whether that is achievable by relaxing the bounds, or
requires merging the degenerate families, is the question posed in
Sec.~\ref{sec:bbbar:degeneracy} and answered by
\texttt{scripts/scan\_bbbar\_bounds.py}}

\subsection{Continuum normalisation}
\label{sec:syst:continuum}
\AnalysisTBD{the continuum normalisation constraint in the workspace is
one-sided, $+15\%/-0\%$, whereas the analysis context document describes it
as $\pm 15\%$.  Confirm which is intended and record the source of the
$15\%$}

\subsection{Fake-$D$ normalisation}
\label{sec:syst:faked}
\AnalysisTBD{the analysis context document records a candidate $\pm 3.5\%$
prior from the $D$-mass sideband, but no such constraint exists in the
workspace generator, where fake-$D$ is an unconstrained normalisation.
Confirm whether the sideband channel is intended to determine this
normalisation in situ, in which case the prior is unnecessary}
